\section{Conclusion}
\label{sec:conclusion}

We have presented the first systematic 50-state empirical study of algorithmically
generated state house maps using the minimum-edge-cut recursive bisection
algorithm. The results establish three main findings.

First, the \texttt{redist} pipeline scales gracefully to the state legislative
context. All 50 state house chambers are successfully redistricted, with
population balance within $\pm 0.5\%$ in every state and runtimes under
8 minutes per state at single-seed operation.

Second, algorithmic state house maps are more compact than both congressional
algorithmic maps ($\overline{PP} = 0.381$--$0.401$ vs.\ $0.361$) and enacted
state house maps ($\overline{PP} = 0.381$ vs.\ $0.300$ for the 38-state
comparison). This counterintuitive result is explained by scale effects (F.5):
smaller districts have relatively shorter perimeters when drawn from compact
sub-regions of the state, and recursive bisection naturally exploits this by
hierarchically partitioning the state from large to small geographic scales.

Third, the algorithmic maps split counties 18\% fewer times on average than
enacted maps, consistent with the algorithm's lack of partisan incentives to
crack geographic communities across administrative boundaries.

These results position the F track as a credible algorithmic baseline for
state legislative redistricting reform. A state redistricting commission
seeking a compact, population-balanced, county-respecting benchmark plan
can generate one using \texttt{redist state --chamber house --resolution block\_group}
in under 10 minutes per state.

\paragraph{Future work.}
Immediate extensions include: partisan outcome analysis using precinct-level
data (requires the political data analysis pipeline, out of scope here);
VRA compliance assessment (F.6); bicameral nesting (F.2); and resolution
refinement for the five states where block-group resolution is inadequate (F.3).
A longer-term goal is a multi-seed ensemble analysis for the high-$k$ chambers,
building on the ensemble methodology of the G track
\citep{dellaLibera2026ensemble}.
